\textbf{Objective.} To measure whether multimodal large language models (MLLMs) that read a 12-lead ECG correctly abandon that reading after a clinician's wrong suggestion, and whether this depends on hint style, diagnosis, model and language (Thai vs English).

\textbf{Methods.} We rendered 101 PTB-XL ECGs in eight groups, including 25 acute-stage and 10 old infarctions, as paper images. Eight MLLMs read each case three times. Cases read correctly in all three runs received a wrong hint (three styles) or a correct hint, in Thai and in English. The primary outcome was the share of responses switching to the hinted diagnosis. A sub-study used open-ended questions with casual remarks.

\textbf{Results.} Baseline accuracy was 3--62\%. Among baseline-correct reads, the five current models switched to the wrong hint in 0.9--6.6\% of wrong-hint responses, close to their no-hint drift. Two previous-generation models switched in 25\% and 32\%. Anecdotal hints (odds ratio 1.9) and Thai hints (1.7) raised the odds of deference. Deference concentrated on acuity: STEMI hints converted old infarctions far more often than normal tracings, and one model withdrew the STEMI call on 19\% of acute-stage infarctions after a ``normal'' hint. One model was contrarian, abandoning correct STEMI reads in 53 of 96 responses when a correct hint agreed. In open-ended use, accuracy fell and casual remarks did not change conclusions.

\textbf{Conclusions.} Current leading MLLMs rarely abandon a correct ECG read under a clinician's wrong suggestion, but deference persists in older models, rises with anecdotal and non-English hints, and targets the acuity judgement that decides reperfusion.
